When an employee concludes that the organization has broken an ethical commitment, the concern can travel in several directions. It can go to a supervisor or an ethics office, where it may be fixed quietly. It can go to the public, through social media or the press, where it may damage the organization's reputation but also force change. It can leave with the employee, or it can stay unspoken. Which of these happens matters to the organization, which learns about problems only through the first route, and to the employee, who faces different risks on each. Hirschman's classic analysis made a related point about organizations in decline: members who speak tell the organization what has gone wrong, while members who leave signal only that something has \citep{hirschman1970exit}. Research on voice and silence has shown that employees often keep concerns to themselves because speaking up feels risky or pointless \citep{morrison2000organizational,milliken2003exploratory}, and research on whistleblowing has shown that people who disclose outside the organization differ from those who report inside it in their circumstances and in the retaliation they suffer \citep{dworkin1998internal}. The direction an ethical complaint takes is a workplace ethics question in its own right.

Discussions of Generation Z at work often start from a different question: whether this cohort cares more about ethics than earlier ones. The evidence gives that premise little support. A large time-lag study did not find that a recent cohort placed more weight on altruistic work values than its predecessors \citep{twenge2010generational}, meta-analytic work has found generational differences in work attitudes to be small \citep{costanza2012generational}, and critical reviews argue that differences attributed to generations are hard to separate from age, period, and career stage and are often better explained by those factors \citep{lyons2014generational,rudolph2021generations}. Building a model of Generation Z's workplace ethics on an assumed difference in values would rest on weak ground.

Two terms need to be pinned down before going further. By workplace ethics we mean the standards of right conduct that govern how an organization and its members treat employees, customers, and others affected by their work, together with the commitments an organization makes about honoring those standards. Research on behavioral ethics has studied how individual, group, and organizational factors shape conduct against such standards \citep{trevino2006behavioral}; our focus is narrower, on what employees do when they believe the organization itself has fallen short of a commitment it made. By Generation Z we mean the cohort that followed the Millennials into the labor market. The birth years used to mark its boundaries differ from source to source, and critics note that such boundaries are set by convention rather than derived from theory \citep{rudolph2021generations}. We therefore use the label descriptively and make no assumption that it marks a group with shared psychological traits. Whether it marks anything that matters for ethical voice is one of the questions the model is built to test.

A more tractable question concerns venue rather than values. A UK practitioner survey reported that its youngest respondents were less likely than older ones to say they would raise concerns with their employer, yet more likely to say they would post about workplace problems on social media or contact journalists \citep{protect2025talkin}. The survey was not peer reviewed and measured intentions at one point in time, so it cannot tell us whether the pattern reflects cohort, age, or the simple fact that young workers tend to be newcomers in junior roles. It does describe a tension worth explaining: reluctance to speak inside combined with readiness to speak outside. Existing theory offers pieces of an explanation. Psychological contract research explains how broken ethical commitments are experienced \citep{thompson2003violations}, and the exit-voice-loyalty tradition explains how people choose among responses to organizational decline \citep{hirschman1970exit,turnley1999impact}. Neither addresses why an employee would choose a public venue over an internal one.

This paper develops a model to address that question. Our central argument is that the venue in which an ethical promise was made shapes the venue in which its breach is contested. Organizations promote what they stand for as employers to outside audiences through employer branding \citep{backhaus2004conceptualizing}, and job seekers respond to signals about corporate social performance \citep{jones2014why}. We assume, without claiming it has been shown, that for many recruits an employer's ethical commitments are therefore first met in public form. We introduce promise publicity to capture this and propose that breaches of publicly made ethical promises are more likely to be felt as violations and more likely to be voiced publicly, a relationship we call venue congruence. Internal voice safety sets the boundary on this effect. Where neither venue seems viable, we expect exit or silence. Finally, rather than assuming that Generation Z behaves differently, we state two rival propositions, one attributing the pattern to cohort and one to career stage, and outline how research could decide between them.

The paper makes three contributions, each of which remains to be tested.

\begin{itemize}
\item It adds promise publicity to psychological contract theory as a feature of how ideological contracts are formed. Prior work has established that contracts form partly from pre-entry schemas and recruitment messages \citep{rousseau2001schema} and that ideological breach triggers compensatory or corrective responses \citep{yang2022going,deng2023serving}; we propose that the public or private origin of the ethical promise affects both the intensity of violation and the venue of the response.
\item It extends exit-voice-loyalty reasoning, which organizational research has applied mainly to internal voice \citep{morrison2014employee}, by treating the choice between internal and public ethical voice as the outcome to be explained, with internal voice safety as its boundary and exit and silence as the alternatives when neither venue appears viable. Studies of ethical culture have shown that culture relates to internal versus external responses \citep{kaptein2011inaction}; our model adds the promise's origin as a second influence on that choice.
\item It recasts the Generation Z question as a contest between rival propositions, following critics who argue that generational labels are often confounded with age and career stage \citep{rudolph2017considering,parry2011generational}, and it differs from the closest existing study of online whistleblowing and ideological contracts \citep{yang2026managing} by modeling the internal alternative, the exit and silence outlets, and the cohort question.
\end{itemize}
